\section{Related work}
\label{sec:related-work}

\subsection{Landscape representations and automated algorithm selection}

Fitness-landscape analysis describes properties of an optimization problem that
may be relevant to search, including ruggedness, modality, neutrality, fitness
distributions, and variable interactions.  The corresponding measures differ
in their assumptions, sampling procedures, dependence on a search process, and
computational demands \cite{malanSurveyTechniquesCharacterising2013}.  ELA
operationalizes this idea through numerical descriptors computed from sampled
objective values.  Its original formulation proposed relatively inexpensive,
computer-generated features as a route from problem characterization to
automatic algorithm selection \cite{mersmannExploratoryLandscapeAnalysis2011}.
Software such as \texttt{pflacco} now exposes multiple feature families with
different sample requirements and domains of applicability
\cite{pragerPflaccoFeaturebasedLandscape2024}.  Estimated ELA values can also
depend materially on the sampling strategy, and a mismatch between the designs
used for training and deployment can degrade the resulting predictor
\cite{renauExploratoryLandscapeAnalysis2020}.  The sample design is therefore
part of the representation, not a neutral implementation detail.

In automated algorithm selection (AAS), a conventional pipeline computes an
instance representation and predicts which portfolio member should be used.
ELA combined with supervised learning has been evaluated against the single
best solver (SBS) and virtual best solver (VBS), often under leave-one-function-
out or related problem-oriented splits
\cite{kerschkeAutomatedAlgorithmSelection2019}.  More recent systems retain the
sample--feature--selector structure while changing the portfolio, feature set,
or learning model \cite{guoAutomatedAlgorithmSelection2025}.  These studies
establish the downstream use of landscape information.  They do not determine
whether an independently sampled descriptor query should be acquired at a
particular state before its features exist.

\subsection{Feature acquisition as a resource cost}

Landscape information is acquired rather than supplied.  Objective evaluations
used for feature sampling reduce the budget left for optimization, while
feature computation adds computational demand.  Reusing the sample
as an optimizer population can reduce the effective FE charge in a particular
pipeline \cite{kerschkeAutomatedAlgorithmSelection2019}, but this does not make
the representation cost-free in general.  Runtime and peak memory are therefore
measured separately in the present study.  A recent preprint study of feature
budgets found
that the preferred sampling fraction varied with the portfolio, problem set,
dimension, target budget, and selection scenario
\cite{blomInfluenceFeatureComputation2026}.  Similarly, ELA-based selection
experiments may restrict feature groups when the acquisition budget is limited
\cite{guoAutomatedAlgorithmSelection2025}.

These findings motivate cost-aware evaluation, but their decision variable is
usually the portfolio action after features have been obtained.  In contrast,
Decision-before-Feature compares executing and skipping the same prespecified
query from a shared search state.  Query FEs enter through the reduced
continuation budget rather than a second penalty; normalized time is penalized
explicitly, while memory is reported and examined only in sensitivity analysis.
The target is therefore not a classifier for whether a feature family is
``cheap.''

\subsection{Trajectory-based and per-run algorithm selection}

The closest antecedents replace an independent landscape sample with
information collected during an initial optimizer run.  Jankovic et al. begin
with CMA-ES, compute trajectory-based landscape features, predict the
performance of candidate second-stage optimizers, and warm-start the selected
one \cite{jankovicTrajectorybasedAlgorithmSelection2022}.  In that formulation,
the trajectory representation is constructed for every run at a fixed switch
point; the decision is which second-stage optimizer to use, not whether to
construct the representation.  Kostovska et al.
extend this per-run setting to compare trajectory ELA, time series of CMA-ES
internal state variables, and their combination
\cite{kostovskaPerrunAlgorithmSelection2022}.  The latter study also reports
that its BBOB findings did not transfer directly to YABBOB, showing that such
transfer cannot be assumed; trajectory-coverage mismatch is discussed there as
a possible explanation.  Both internal evaluations retain the same BBOB
functions across training and test partitions: the CEC study groups
five-dimensional observations by instance identifier, while the PPSN study uses
run- and instance-level partitions in five and ten dimensions.  These protocols
answer different questions from the present function-family OOF design; they are
antecedents for representation and warm-starting, not precedents for the present
evaluation design.

Probing trajectories provide another algorithm-centric representation.  A
short run records an objective-value sequence, and the resulting time series
supports a later portfolio choice
\cite{renauUtilityProbingTrajectories2024}.
A subsequent benchmark compared multiple time-series classifiers under both
leave-one-instance-out and the harder leave-one-problem-out evaluation
\cite{renauAlgorithmSelectionProbing2025b}.  This result is relevant here
because it shows that the predictive model and problem partition remain part
of a trajectory selector's empirical specification.  Nevertheless, probing is
performed before every selection, the target is the best portfolio algorithm,
and the output is an optimizer choice.  The probe itself still consumes
evaluations even though it is not a separate landscape sample.  It does not construct paired
skip/query utility for a separate landscape query or decide whether that query
should be acquired.

\subsection{Trajectory representations and behavior characterization}

Other work studies what optimizer histories represent rather than which action
to take.  Search trajectory networks summarize visited regions and transitions
for comparative analysis and visualization \cite{ochoaSearchTrajectoryNetworks2021}.
DynamoRep concatenates generation-wise coordinate and fitness summaries, thereby
preserving longitudinal population dynamics without a separate objective sample;
Opt2Vec instead learns a representation from populations observed during a run
\cite{cenikjDynamoRepTrajectoryBasedPopulation2023,korosecOpt2VecContinuousOptimization2024}.
DynamoRep fits separate classifiers for each optimizer and evaluates
three-dimensional BBOB class identification, while Opt2Vec likewise reports
problem classification.  These tasks neither establish a fixed-dimensional,
algorithm-agnostic state representation nor predict the utility of acquiring a
separate landscape query.

Behavioral analysis provides a complementary vocabulary.  Hayward and
Engelbrecht characterize exploration, exploitation, locality, communication,
and evaluation effort with a suite of whole-run measures, explicitly treating
behavior as problem dependent \cite{haywardDeterminingMetaheuristicSimilarity2025}.
Some of those measures rely on a known optimum, persistent individual identity,
or algorithm interaction structure and hence are unsuitable as pre-query inputs
here.  Mbasso et al. use distance-to-reference decay, directional entropy, and
stagnation to diagnose late-stage behavior and trigger optimizer-specific
interventions \cite{mbassoHowMetaheuristicsExploit2026}.  Only the distance-decay
and stagnation semantics motivate the present state representation: distance is
redefined relative to the current population best, while particle-wise
directional entropy is excluded because it requires persistent cross-generation
identities.  The reported interventions do not validate the study-specific
Search Maturity construction or establish that behavior predicts landscape-query
utility.  More generally, computability alone does not
make a behavior indicator informative or complementary: sensitivity,
redundancy, single-run computability, and the information available at
deployment all constrain its use
\cite{haywardSurveyAnalysisMetaheuristic2026}.

\subsection{Adaptive search and dynamic algorithm selection}

Landscape-aware adaptive methods change the search mechanism after observing
landscape or performance information.  Examples include mutation-operator
selection in differential evolution and adaptation of subpopulation number and
partitioning in artificial bee colony optimization
\cite{sallamLandscapebasedAdaptiveOperator2017,zhouAdaptiveMultipopulationArtificial2024}.
Adaptive landscape analysis instead studies how locally sampled landscape
descriptors evolve during an optimizer run, with the longer-term aim of online
selection or configuration \cite{jankovicAdaptiveLandscapeAnalysis2019a}.  It
therefore computes descriptors on 2,000 additional points sampled from the
current CMA-ES distribution at prescribed target levels.  ``Cheap'' denotes
that the chosen descriptor families need no further evaluations beyond that
sample, not that acquisition is free.  The preliminary study neither makes a
pre-descriptor acquisition decision nor couples the descriptors to a paired
skip/query utility.
Dynamic algorithm selection can make repeated choices during a run; RL-DAS, for
example, represents the process as a reinforcement-learning problem over
differential-evolution variants using landscape and algorithm-history
information \cite{guoDeepReinforcementLearning2024b}.  Dynamic algorithm
configuration makes the control problem explicit through a state space, an
action space over hyperparameter settings or components, a reward, and a
decision frequency; DACBench includes both CMA-ES step-size control and repeated
selection among modular evolutionary-algorithm components
\cite{eimerDACBenchBenchmarkLibrary2021}.  Parameter control and MetaBBO cover
still broader choices over parameters, components, structures, and learned
update rules
\cite{gomespereiradelacerdaSystematicLiteratureReview2021,yangMetablackboxOptimizationEvolutionary2025}.
The broader adaptive-operator-selection literature likewise separates
stateless credit-based rules from state-based policies that map solution,
problem, or optimization-process features to an operator
\cite{peiAdaptiveOperatorSelection2025}.  In both cases, the selected action is
a search operator; the acquisition of an independent descriptor query is not
the action being evaluated.

Those approaches optimize a search action or solver mechanism.  The present
method uses offline supervised learning and leaves the portfolio algorithms
unchanged.  Its gate decides whether to acquire information; only after a call
does a separate, fixed action-loss selector choose between native continuation
and population-transfer initialization.  Thus, its action is neither an
optimizer hyperparameter nor a repeated component choice.  Table~\ref{tab:related-comparison}
summarizes this distinction for the closest task formulations.

\begin{table*}[t]
\caption{Closest task formulations compared by decision object, information
time, acquisition cost, and state transition. ``Extra FEs'' denotes an
independent landscape sample beyond an existing prefix or probe.}
\label{tab:related-comparison}
\centering
\footnotesize
\begin{tabularx}{\textwidth}{@{}l Y Y c Y@{}}
\toprule
Approach & Decision target & Information at decision & Extra FEs
& Frequency and transition \\
\midrule
Static ELA-based AAS
\cite{kerschkeAutomatedAlgorithmSelection2019,guoAutomatedAlgorithmSelection2025}
& Portfolio algorithm
& Acquired landscape descriptors
& Usually yes
& Once, before the chosen optimizer starts \\
Jankovic et al. \cite{jankovicTrajectorybasedAlgorithmSelection2022}
& Second-stage algorithm
& ELA computed from a fixed CMA-ES prefix
& No separate sample
& Once, with algorithm-specific warm-starting \\
Kostovska et al. \cite{kostovskaPerrunAlgorithmSelection2022}
& Second-stage algorithm
& Trajectory ELA and/or CMA-ES state series
& No separate sample
& Once, with algorithm-specific warm-starting \\
Probing trajectories \cite{renauUtilityProbingTrajectories2024,renauAlgorithmSelectionProbing2025b}
& Portfolio algorithm
& Objective-value sequence from a short probe
& No separate landscape sample
& Once, after probing \\
RL-DAS \cite{guoDeepReinforcementLearning2024b}
& Differential-evolution variant
& Landscape, search history, and context
& Not isolated as a fixed query
& Repeated online scheduling \\
DACBench \cite{eimerDACBenchBenchmarkLibrary2021}
& Hyperparameter or algorithm component
& Target-algorithm state and instance context
& Not a landscape-query decision
& Repeated dynamic configuration \\
Decision-before-Feature
& Execute or skip one fixed query
& Pre-query algorithm-agnostic behavior
& Only when called
& Opportunities until one call; native continuation or population transfer \\
\bottomrule
\end{tabularx}
\end{table*}

\subsection{Evaluation splits and the remaining gap}

Representation quality and selector performance depend strongly on how
problems are divided between training and evaluation.  In an affine-BBOB study,
several landscape representations that performed well under easier splits did
not outperform SBS under the most demanding problem-oriented splits
\cite{cenikjLandscapeFeaturesSingleobjective2025}.  A recent representation
survey likewise highlights inconsistent sampling and split protocols, heavy
dependence on BBOB, and limited cross-benchmark evidence
\cite{cenikjSurveyFeaturesUsed2026}.  These observations motivate function-
family out-of-fold selection and a held-out evaluation; they do not guarantee
that the resulting gate transfers.
Benchmarking work further shows that leave-one-instance-out evaluation can
reward representations that recover a problem class shared by related training
and test instances without learning to rank algorithms on unseen problems, and
that scale-sensitive prediction errors need not measure selection quality
\cite{petelinPitfallsBenchmarkingAlgorithm2025}.  The present study therefore
uses function-family partitions and selects the Decision Model by policy utility
rather than raw action-loss regression error.  This design addresses those two
failure modes without establishing generalization in advance.

Among the literature reviewed here, we did not identify a method that makes the
same decision at the same information time: before descriptors from a
prespecified independent query exist, use only algorithm-agnostic behavior to
predict whether the query's downstream performance difference justifies its FE
and computational cost.  Decision-before-Feature addresses this corpus-scoped
gap with matched-state offline labels, an explicit separation between the query
gate and the downstream selector, a time-only comparator, and nested
function-family evaluation.  The literature supports this formulation and its
design constraints; it does not supply evidence for the pending RQ1--RQ5
outcomes or for held-out and cross-benchmark generalization.
